\documentclass[twocolumn,aps,pra,floatfix,amsmath,amssymb,longbibliography]{revtex4-2}

\usepackage[T1]{fontenc}
\usepackage[
  breaklinks=true,
  colorlinks=true,
  citecolor=black,
  linkcolor=black,
  urlcolor=black
]{hyperref}

\newcommand{\ii}{\mathrm{i}}
\newcommand{\C}{\mathbb{C}}
\newcommand{\R}{\mathbb{R}}

\begin{document}

\appendix

\section{The opening paragraph worked out}
\label{app:spin1}

This appendix derives each statement of the first paragraph of Sec.~\ref{sec:intro}.  We write $\langle\phi|\psi\rangle=\phi^\dagger\psi$; $\varepsilon_{kab}$ is the Levi-Civita symbol on $\{x,y,z\}$ with $\varepsilon_{xyz}=1$; and dot and cross products of real vectors are extended linearly to complex vectors.

In brief, the Kochen--Specker realization can only produce real outcome vectors.  In the Cartesian basis, the three outcome states of every squared-spin measurement are real (up to phases), so this scheme cannot realize measurement arrangements that require complex vectors.  The restriction comes from the squared-spin observables, not from spin~1 itself: a nondegenerate observable that is a real quadratic form in the spin components has a unique eigenbasis, which is real up to phases, whereas $S_{\mathbf n}$ itself, which is linear in $\mathbf S$, has complex eigenstates for the eigenvalues $\pm\hbar$.  ``Real'' refers to the Cartesian basis; in the spherical basis the same outcome states generally have complex coordinates.  What does not depend on the basis is whether there is a basis in which all outcome states of an arrangement are real (up to phases); see the last subsection.

\subsection{Cartesian basis}
In the \emph{spherical} ($S_z$) eigenbasis $(|{+1}\rangle,|0\rangle,|{-1}\rangle)$ with Condon--Shortley phases, the spin-1 matrices are~\cite{Sakurai-1994}
\begin{equation}
\setlength{\arraycolsep}{2.5pt}
\begin{aligned}
 S_x^{\mathrm{sph}}&=\frac{\hbar}{\sqrt{2}}\!\begin{pmatrix}0&1&0\\1&0&1\\0&1&0\end{pmatrix}\!,&
 S_y^{\mathrm{sph}}&=\frac{\hbar}{\sqrt{2}}\!\begin{pmatrix}0&-\ii&0\\\ii&0&-\ii\\0&\ii&0\end{pmatrix}\!,\\
 S_z^{\mathrm{sph}}&=\hbar\!\begin{pmatrix}1&0&0\\0&0&0\\0&0&-1\end{pmatrix}\!.
\end{aligned}
\label{eq:sph-matrices}
\end{equation}
Solving $S_k|k\rangle=0$ for $k=x,y,z$ gives the normalized states
\begin{equation}
 |x\rangle=\frac{|{-1}\rangle-|{+1}\rangle}{\sqrt2},\;\;
 |y\rangle=\ii\,\frac{|{+1}\rangle+|{-1}\rangle}{\sqrt2},\;\;
 |z\rangle=|0\rangle,
 \label{eq:cart-from-sph}
\end{equation}
which are orthonormal.  Let $U$ be the unitary matrix whose columns are the coordinates of $|x\rangle,|y\rangle,|z\rangle$ in this spherical basis, Eq.~\eqref{eq:cart-from-sph}.  Direct multiplication, $S_k=U^\dagger S_k^{\mathrm{sph}}U$, gives the spin-1 matrices in the \emph{Cartesian basis} $\{|x\rangle,|y\rangle,|z\rangle\}$:
\begin{equation}
\setlength{\arraycolsep}{1.8pt}
 S_x=\hbar\!\begin{pmatrix}0&0&0\\0&0&-\ii\\0&\ii&0\end{pmatrix}\!,\;
 S_y=\hbar\!\begin{pmatrix}0&0&\ii\\0&0&0\\-\ii&0&0\end{pmatrix}\!,\;
 S_z=\hbar\!\begin{pmatrix}0&-\ii&0\\\ii&0&0\\0&0&0\end{pmatrix}\!,
 \label{eq:cart-matrices}
\end{equation}
that is, $(S_k)_{ab}=\langle a|S_k|b\rangle=-\ii\hbar\,\varepsilon_{kab}$.  These matrices are purely imaginary and antisymmetric, $S_k^T=-S_k$.  Rephasing $|a\rangle\mapsto e^{\ii\phi_a}|a\rangle$ multiplies $(S_k)_{ab}$ by $e^{\ii(\phi_b-\phi_a)}$, so Eq.~\eqref{eq:cart-matrices} fixes the phases in Eq.~\eqref{eq:cart-from-sph} up to one overall phase.  From here on we work in this Cartesian basis, so $|x\rangle,|y\rangle,|z\rangle$ form the standard coordinate basis $(1,0,0)^T$, $(0,1,0)^T$, $(0,0,1)^T$ of $\C^3$ (``Cartesian columns''), whereas the spherical basis states $(|{+1}\rangle,|0\rangle,|{-1}\rangle)$ are the columns of $U^\dagger$.

\subsection{Squared spin components}
For a real unit vector $\mathbf n\in\R^3$, Eq.~\eqref{eq:cart-matrices} gives
\begin{equation}
 S_{\mathbf n}:=\mathbf n\cdot\mathbf S
 =\ii\hbar\!\begin{pmatrix}
  0 & -n_z & n_y \\
  n_z & 0 & -n_x \\
  -n_y & n_x & 0
 \end{pmatrix}\!,
 \qquad
 S_{\mathbf n}\psi=\ii\hbar\,\mathbf n\times\psi
 \label{eq:Sn}
\end{equation}
for every Cartesian column $\psi\in\C^3$.  Applying $S_{\mathbf n}$ twice,
\begin{equation}
 S_{\mathbf n}^2\psi=-\hbar^2\,\mathbf n\times(\mathbf n\times\psi)
 =\hbar^2\big(\psi-\mathbf n\,(\mathbf n\cdot\psi)\big),
 \label{eq:Sn2-action}
\end{equation}
so $S_{\mathbf n}^2=\hbar^2(I_3-P_{\mathbf n})$ with $P_{\mathbf n}:=\mathbf n\mathbf n^T$, the projector onto the real vector $\mathbf n$.  For two real unit vectors,
\begin{equation}
 [S_{\mathbf n_1}^2,S_{\mathbf n_2}^2]
 =\hbar^4[P_{\mathbf n_1},P_{\mathbf n_2}]
 =\hbar^4(\mathbf n_1\cdot\mathbf n_2)\big(\mathbf n_1\mathbf n_2^T-\mathbf n_2\mathbf n_1^T\big).
 \label{eq:comm-Sn2}
\end{equation}
For $\mathbf n_1\neq\pm\mathbf n_2$ the last factor is nonzero (applied to $\mathbf n_2$ it gives $\mathbf n_1-(\mathbf n_1\cdot\mathbf n_2)\mathbf n_2\neq\mathbf 0$).  So $S_{\mathbf n_1}^2$ and $S_{\mathbf n_2}^2$ commute if and only if $\mathbf n_1\perp\mathbf n_2$, and three distinct squared components $S_{\mathbf e_j}^2$ commute if and only if $\{\mathbf e_1,\mathbf e_2,\mathbf e_3\}$ is an orthonormal triad.  For such a triad $P_{\mathbf e_1}+P_{\mathbf e_2}+P_{\mathbf e_3}=I_3$, hence
\begin{equation}
 S_{\mathbf e_1}^2+S_{\mathbf e_2}^2+S_{\mathbf e_3}^2=2\hbar^2I_3,
 \qquad
 S_{\mathbf e_j}^2\,\mathbf e_k=\hbar^2(1-\delta_{jk})\,\mathbf e_k.
 \label{eq:triad}
\end{equation}
Each $S_{\mathbf e_j}^2$ takes the values $0$ or $\hbar^2$, and exactly one of the three is $0$.  Kochen and Specker measure all three at once through the Hamiltonian $H=aS_{\mathbf e_1}^2+bS_{\mathbf e_2}^2+cS_{\mathbf e_3}^2$ with distinct real $a,b,c$~\cite{kochen1}, whose nondegenerate eigenvalues $\hbar^2(b+c)$, $\hbar^2(c+a)$, $\hbar^2(a+b)$ belong to $\mathbf e_1$, $\mathbf e_2$, $\mathbf e_3$.  The squared-spin measurement along the triad is therefore the three-outcome projective measurement in the basis $\{\mathbf e_1,\mathbf e_2,\mathbf e_3\}$, outcome $j$ meaning $S_{\mathbf e_j}^2=0$.  More generally, each $S_k$ is purely imaginary and antisymmetric, so every product $S_jS_k$ is real, with $(S_jS_k)^T=S_kS_j$.  Hence a real quadratic form $\sum_{j,k}q_{jk}S_jS_k$ with symmetric $(q_{jk})$ is a real symmetric matrix, and if it is nondegenerate, its eigenbasis is real up to phases.

\subsection{Zero-spin states are real}
Let $|\mathbf n\rangle:=n_x|x\rangle+n_y|y\rangle+n_z|z\rangle$, whose Cartesian column is $\mathbf n$ itself.  By Eq.~\eqref{eq:Sn}, $S_{\mathbf n}\mathbf n=\ii\hbar\,\mathbf n\times\mathbf n=\mathbf 0$.  Conversely, write a normalized $\psi$ as $\psi=\mathbf x+\ii\mathbf y$ with $\mathbf x,\mathbf y\in\R^3$.  Then $S_{\mathbf n}\psi=\mathbf 0$ means $\mathbf n\times\mathbf x=\mathbf n\times\mathbf y=\mathbf 0$, so $\psi=(\alpha+\ii\beta)\mathbf n$ with $\alpha^2+\beta^2=1$: the zero-spin state along $\mathbf n$ is $|\mathbf n\rangle$, unique up to a phase.  Its mean spin vanishes, because $\langle\mathbf n|S_k|\mathbf n\rangle=\mathbf n^TS_k\mathbf n$ is a number equal to its own transpose $\mathbf n^TS_k^T\mathbf n=-\mathbf n^TS_k\mathbf n$; for instance, $\langle\mathbf n|S_x|\mathbf n\rangle=\ii\hbar(n_zn_y-n_yn_z)=0$.

\subsection{Squared-spin measurements as real bases}
By the two preceding subsections, the outcome states of the squared-spin measurement along $\{\mathbf e_1,\mathbf e_2,\mathbf e_3\}$ are the real vectors $\mathbf e_1,\mathbf e_2,\mathbf e_3$, and every real orthonormal basis of $\R^3$ arises in this way.  For real unit vectors, $\langle\mathbf n_1|\mathbf n_2\rangle=\mathbf n_1\cdot\mathbf n_2$, so $\mathbf n_1,\mathbf n_2$ give the same ray if and only if $\mathbf n_1=\pm\mathbf n_2$, and orthogonal rays if and only if $\mathbf n_1\cdot\mathbf n_2=0$.  Shared outcomes and orthogonalities among squared-spin measurements are therefore exactly those among the corresponding real bases: a finite family of squared-spin measurements is a finite family of real orthonormal bases of $\R^3$.

\subsection{Measurement arrangements}
A finite arrangement of three-outcome measurements is a hypergraph whose vertices are outcomes and whose three-element edges are the measurements (contexts).  A faithful realization in $\mathbb F^3$ ($\mathbb F=\R$ or $\C$) is a faithful orthogonal representation in the sense of Definition~\ref{def:for}: it assigns unit vectors $\psi_v\in\mathbb F^3$ to the vertices so that distinct vertices give distinct rays, and two vertices give orthogonal vectors if and only if they belong to a common context.  By the previous subsection, an arrangement can be realized with squared-spin measurements alone if and only if it has a faithful realization in $\R^3$.

\subsection{Mean spin}
For normalized $\psi=\mathbf x+\ii\mathbf y$, the matrix elements $(S_k)_{ab}=-\ii\hbar\,\varepsilon_{kab}$ give $\langle\psi|S_k|\psi\rangle=-\ii\hbar\sum_{a,b}\varepsilon_{kab}\,\overline{\psi_a}\,\psi_b$, that is~\cite{ZiboldEtAl2016SpinNematic},
\begin{equation}
 \langle\psi|\mathbf S|\psi\rangle
 =-\ii\hbar\,\psi^*\times\psi
 =2\hbar\,\mathbf x\times\mathbf y,
 \label{eq:mean-spin}
\end{equation}
since $\psi^*\times\psi=(\mathbf x-\ii\mathbf y)\times(\mathbf x+\ii\mathbf y)=2\ii\,\mathbf x\times\mathbf y$.  Thus $\langle\mathbf S\rangle=\mathbf 0$ if and only if $\mathbf x$ and $\mathbf y$ are linearly dependent, that is, if and only if $\psi=e^{\ii\phi}\mathbf n$ for a real unit vector $\mathbf n$.  By contrast, the other two eigenstates of $S_{\mathbf n}$, $\mp(\mathbf e_1\pm\ii\mathbf e_2)/\sqrt2$ with real unit vectors $\mathbf e_1,\mathbf e_2$ satisfying $\mathbf e_1\times\mathbf e_2=\mathbf n$, have $\langle\mathbf S\rangle=\pm\hbar\mathbf n$.  So, in the Cartesian basis, the outcome states of a realization all have zero mean spin if and only if they are all real up to phases.

\subsection{Choice of basis}
Any orthonormal basis $\{W|x\rangle,W|y\rangle,W|z\rangle\}$, with $W$ unitary, can be taken as the Cartesian basis: it is the Cartesian basis of the spin operators $\mathbf S^{(W)}:=W\mathbf SW^\dagger$.  Every spin-1 operator triple on $\C^3$ has this form, because the irreducible spin-1 representation of $\mathfrak{su}(2)$ is unique up to unitary equivalence~\cite{Sakurai-1994,hall-2015}.  With $\psi'=W^\dagger\psi$,
\begin{equation}
 \langle\psi|\mathbf S^{(W)}|\psi\rangle
 =\langle\psi'|\mathbf S|\psi'\rangle
 =2\hbar\,\operatorname{Re}\psi'\times\operatorname{Im}\psi'.
 \label{eq:SW-mean}
\end{equation}
Suppose a faithful realization $\{\psi_v\}$ in $\C^3$ and some $W$ give $\langle\psi_v|\mathbf S^{(W)}|\psi_v\rangle=\mathbf 0$ for all $v$.  Then $W^\dagger\psi_v=e^{\ii\phi_v}\mathbf n_v$ with real unit vectors $\mathbf n_v$, and $|\mathbf n_u\cdot\mathbf n_v|=|\langle\psi_u|\psi_v\rangle|$ for all $u,v$, so the $\mathbf n_v$ form a faithful realization in $\R^3$.  Conversely, a faithful realization in $\R^3$ has zero mean spin at every vertex in the Cartesian basis.  Hence an arrangement has no faithful realization in $\R^3$ if and only if, in every faithful realization in $\C^3$ and whatever basis is taken as Cartesian, at least one outcome state has $\langle\mathbf S^{(W)}\rangle\neq\mathbf 0$.  Because $\operatorname{Tr}\mathbf S^{(W)}=\mathbf 0$, the mean spins of the three outcomes of any context sum to zero, so such a state has a partner with nonzero mean spin in each of its contexts.  To show that the second alternative occurs, it therefore suffices to exhibit a finite hypergraph with a faithful realization in $\C^3$ but none in $\R^3$, as done in the main text.

\bibliography{svozil}

\end{document}
